\section{Instantiation}\label{sec:tsps-instantiation}

\begin{outline}
  \item we instantiate with $\AGHO$ and $\KPW$
  \item $\AGHO$: smallest possible signatures, secret key only field elements, so
the threshold version only needs multiplication of shared field elements
  \item $\KPW$: standard assumption, but secret key contains group elements,
hence sharings of group elements and $\algo{ExpGrp}$
  \item for each: show signing is decomposable, give the threshold scheme in full
  \item reminder: raising a public group element to a shared field element is
each signer raising it to its own share, giving a sharing of the result
\end{outline}

\subsection{Instantiating $\TSPS$ with $\AGHO$}\label{sec:tsps-abe}

\begin{outline}
  \item $\AGHO$ recalled in the Preliminaries (cite C:AGHO11); with $\rho =
1/\omega$ for its signing randomness, a signature on $\vec{\msg}$ is
      %
      \begin{equation*}
      %
	(\widetilde{r}, \widetilde{s}, z) = \left(\ggen^{1/\rho}, \
\ggen^{\tau/\rho}, \ \gen^{\rho} \prod_{i=1}^{\ell} \msg_i^{-\rho\mu_i}\right)
      %
      \end{equation*}
  \item decomposable with $c_0 = (\widetilde{r}, \widetilde{s})$, $c_1 =
\gen^{\rho}$, $e_{1,i} = -\rho\mu_i$, none depending on the message
  \item everything offline is a field element or a generator raised to one, so
no $\algo{ExpGrp}$
  \item the costly step is $1/\rho$: directly it's a multiplication and an
opening, then $\tau/\rho$ one more multiplication, three rounds; multiplying
$\tau$ by the same mask as $\rho$ in the first round gives $\tau/\rho$ once the
mask is opened, so two rounds
  \item result: threshold $\AGHO$, \Cref{prot:tsps-abe}
\end{outline}

\begin{protocolbox}{The threshold structure-preserving signature scheme from AGHO}{tsps-abe}

  \begin{description}

    \item[$(\pk, \shr{\sk}) \sample \TSPSgen(\secparam, \ell, t, n)$ :] The
signers compute
      %
      \begin{equation*}
      %
	\shr{\mu_1}, \ldots, \shr{\mu_\ell}, \shr{\tau} \sample
\algo{RandShare}(), \quad \shr{v} = \gen^{\shr{\tau}}, \quad
\shr{\widetilde{u}_i} = \ggen^{\shr{\mu_i}} \ \ (i \in [1,\ell]),
      %
      \end{equation*}
      %
      and open $\pk = (\widetilde{u}_1, \ldots, \widetilde{u}_\ell, v)$. Each
signer keeps $\shr{\sk} = (\shr{\mu_1}, \ldots, \shr{\mu_\ell}, \shr{\tau})$.

    \item[$\shr{\state} \sample \TSPSoffsign(\shr{\sk})$ :] The signers sample
$\shr{\rho}, \shr{\beta} \sample \algo{RandShare}()$, compute
      %
      \begin{equation*}
      %
	\shr{\rho\beta}, \ \shr{\rho\mu_i} \ \ (i \in [1,\ell]), \
\shr{\tau\beta} \sample \algo{Mul}
      %
      \end{equation*}
      %
      in one round, and in a second round open $\gamma = \rho\beta$, aborting if
$\gamma = 0$. Each signer then computes locally
      %
      \begin{equation*}
      %
	\shr{1/\rho} = \gamma^{-1}\shr{\beta}, \quad \shr{\tau/\rho} =
\gamma^{-1}\shr{\tau\beta}, \quad \shr{\widetilde{r}} = \ggen^{\shr{1/\rho}},
\quad \shr{\widetilde{s}} = \ggen^{\shr{\tau/\rho}}, \quad \shr{\gen^{\rho}} =
\gen^{\shr{\rho}},
      %
      \end{equation*}
      %
      and stores $\shr{\state} = (\shr{\widetilde{r}}, \shr{\widetilde{s}},
\shr{\gen^{\rho}}, \shr{\rho\mu_1}, \ldots, \shr{\rho\mu_\ell})$.

    \item[$\SPSsig_i \leftarrow \TSPSonsign(\sk_i, \state_i, \vec{\msg})$ :]
Each signer computes
      %
      \begin{equation*}
      %
	\shr{z} = \shr{\gen^{\rho}} \cdot \prod_{i=1}^\ell
\msg_i^{-\shr{\rho\mu_i}},
      %
      \end{equation*}
      %
      and outputs $(\shr{\widetilde{r}}, \shr{\widetilde{s}}, \shr{z})$.

    \item[$\SPSsig \leftarrow \TSPScombine(\{\SPSsig_i\}_{i \in Q})$ :]
Reconstruct $\widetilde{r}$, $\widetilde{s}$ and $z$ from the shares of the $t+1$
signers in $Q$, and output $\SPSsig = (\widetilde{r}, \widetilde{s}, z)$.

    \item[$\{0,1\} \leftarrow \TSPSverify(\pk, \vec{\msg}, \SPSsig)$ :] Output
$1$ if and only if
      %
      \begin{equation*}
      %
	\pairing(v, \widetilde{r}) = \pairing(\gen, \widetilde{s}) \ \wedge\
\pairing(z, \widetilde{r}) \prod_{i=1}^{\ell} \pairing(\msg_i, \widetilde{u}_i)
= \pairing(\gen, \ggen).
      %
      \end{equation*}

  \end{description}

\end{protocolbox}

\subsection{Instantiating $\TSPS$ with $\KPW$}\label{sec:tsps-kiltz}

\Cref{prim:sps-kiltz} recalls the $\KPW$ SPS scheme~\cite{C:KilPanWee15}. We write
$k_{i,j}$ for the entry of $\mat{K}$ in row $i$ and column $j$, and
$\mathsf{T}$ denotes transposition. For a vector of group elements $\vec{h} =
(h_1, \ldots, h_{\ell+1})$ and a matrix $\mat{K} \in \Zp^{(\ell+1) \times 2}$,
we write $\vec{h}^{\,\mat{K}}$ for the pair $(\prod_i h_i^{k_{i,1}}, \prod_i
h_i^{k_{i,2}})$, and $(\gen, \vec{\msg})$ for the vector $(\gen, \msg_1, \ldots,
\msg_\ell)$. For a vector $\vec{a}$ of field
elements we write $\gen^{\vec{a}}$ for $(\gen^{a_1}, \gen^{a_2}, \ldots)$.
Products of vectors are taken entry by entry, and for $\vec{h} \in \group{1}^2$
and $\vec{\widetilde{h}} \in \group{2}^2$ we write $\pairing(\vec{h},
\vec{\widetilde{h}}) = \pairing(h_1, \widetilde{h}_1)\pairing(h_2,
\widetilde{h}_2)$.

\begin{primitivebox}{The $\KPW$ SPS scheme}{sps-kiltz}

  \begin{description}

    \item[$(\pk, \sk) \sample \SPSgen(\secparam, \ell)$ :] Sample $\alpha, \beta
\sample \Zp$, $\mat{K} \sample \Zp^{(\ell+1)\times 2}$ and $\mat{\Gamma},
\mat{\Delta} \sample \Zp^{2\times 2}$. Let $\vec{\alpha} = (1,
\alpha)^{\mathsf{T}}$ and $\vec{\beta} = (1, \beta)^{\mathsf{T}}$, and compute
$\vec{\mu} = \mat{K}\vec{\alpha}$, $\vec{\nu} = \mat{\Gamma}\vec{\alpha}$,
$\vec{\omega} = \mat{\Delta}\vec{\alpha}$, $\vec{\xi} =
\vec{\beta}^{\mathsf{T}}\mat{\Gamma}$, and $\vec{\upsilon} =
\vec{\beta}^{\mathsf{T}}\mat{\Delta}$. Output
      %
      \begin{gather*}
      %
	\sk = (\mat{K}, \vec{b}, \vec{x}, \vec{y}) = (\mat{K}, \gen^{(1,
\beta)}, \gen^{\vec{\xi}}, \gen^{\vec{\upsilon}}), \\
      %
	\pk = (\vec{\widetilde{a}}, \vec{\widetilde{u}}, \vec{\widetilde{v}},
\vec{\widetilde{w}}) = (\ggen^{(1, \alpha)}, \ggen^{\vec{\mu}},
\ggen^{\vec{\nu}}, \ggen^{\vec{\omega}}).
      %
      \end{gather*}

    \item[$\SPSsig \sample \SPSsign(\sk, \vec{\msg})$ :] Sample $\rho, \sigma
\sample \Zp^*$ and output $\SPSsig = (\widetilde{q}, \vec{r}, \vec{s},
\vec{z})$, where
      %
      \begin{equation*}
      %
	\widetilde{q} = \ggen^{\sigma}, \quad \vec{r} = \vec{b}^{\,\rho},
\quad \vec{s} = \vec{b}^{\,\rho\sigma}, \quad \vec{z} = (\gen,
\vec{\msg})^{\mat{K}} \cdot \vec{x}^{\,\rho}\, \vec{y}^{\,\rho\sigma}.
      %
      \end{equation*}

    \item[$\{0,1\} \leftarrow \SPSverify(\pk, \vec{\msg}, \SPSsig)$ :] Output
$1$ if and only if
      %
      \begin{equation*}
      %
	\pairing(\vec{z}, \vec{\widetilde{a}}) = \pairing((\gen,
\vec{\msg}), \vec{\widetilde{u}})\, \pairing(\vec{r},\vec{\widetilde{v}})\,
\pairing(\vec{s}, \vec{\widetilde{w}}) \ \wedge \ \pairing(r_1, \widetilde{q})
= \pairing(s_1, \ggen),
      %
      \end{equation*}
      %
      where $r_1$ and $s_1$ are the first entries of $\vec{r}$ and $\vec{s}$, and
$\pairing((\gen, \vec{\msg}), \vec{\widetilde{u}}) = \pairing(\gen,
\widetilde{u}_1) \prod_{i=1}^{\ell} \pairing(\msg_i, \widetilde{u}_{i+1})$.

  \end{description}

\end{primitivebox}

\begin{outline}
  \item $\KPW$ is EUF-CMA under SXDH (cite C:KilPanWee15)
  \item decomposable: $c_0 = (\widetilde{q}, \vec{r}, \vec{s})$, and entry $j$ of
$\vec{z}$ has $c_j = \gen^{k_{1,j}} x_j^{\rho} y_j^{\rho\sigma}$ and $e_{j,i} =
k_{i+1,j}$; none depend on the message, which only enters as public $\msg_i$
raised to shared entries of $\mat{K}$
  \item our protocol and implementation compute $\gen^{k_{1,j}}$ in the online
phase instead, also local
\end{outline}

Unlike $\AGHO$, the secret key contains the group elements $\vec{b}, \vec{x},
\vec{y} \in \group{1}^2$, and the offline phase raises each of them to a shared
exponent. These are shared group elements raised to shared field elements, which
needs $\algo{ExpGrp}$, and so $\KPW$ needs sharings of group elements.
The signers could instead keep the discrete logarithms $\beta, \vec{\xi},
\vec{\upsilon}$ as shares of field elements, so that $\algo{Mul}$ and linear
operations would suffice, but they would then no longer hold shares of the
secret key of $\KPW$ as it is defined. We keep the key as defined.

\begin{outline}
  \item threshold $\KPW$ in \Cref{prot:tsps-kiltz}
  \item offline phase: one multiplication, then one round of eight calls to
$\algo{ExpGrp}_1$
  \item first entry of $\vec{b}$ is the public $\gen$, so two of the eight could
be local, but the implementation makes all eight
\end{outline}

\begin{protocolbox}{The threshold structure-preserving signature scheme from KPW}{tsps-kiltz}

  \begin{description}

    \item[$(\pk, \shr{\sk}) \sample \TSPSgen(\secparam, \ell, t, n)$ :] The
signers sample $\shr{\alpha}$, $\shr{\beta}$ and every entry of $\shr{\mat{K}}$,
$\shr{\mat{\Gamma}}$ and $\shr{\mat{\Delta}}$ with $\algo{RandShare}$, and
compute $\shr{\vec{\mu}}$, $\shr{\vec{\nu}}$, $\shr{\vec{\omega}}$,
$\shr{\vec{\xi}}$ and $\shr{\vec{\upsilon}}$ with one call to $\algo{Mul}$ per
entry, followed by a local addition, since the first entries of $\vec{\alpha}$ and
$\vec{\beta}$ are the public constant $1$. Each signer then
computes locally
      %
      \begin{gather*}
      %
	\shr{\vec{b}} = \gen^{(1, \shr{\beta})}, \quad \shr{\vec{x}} =
\gen^{\shr{\vec{\xi}}}, \quad \shr{\vec{y}} = \gen^{\shr{\vec{\upsilon}}}, \\
      %
	\shr{\vec{\widetilde{a}}} = \ggen^{(1, \shr{\alpha})}, \quad
\shr{\vec{\widetilde{u}}} = \ggen^{\shr{\vec{\mu}}}, \quad
\shr{\vec{\widetilde{v}}} = \ggen^{\shr{\vec{\nu}}}, \quad
\shr{\vec{\widetilde{w}}} = \ggen^{\shr{\vec{\omega}}},
      %
      \end{gather*}
      %
      and the signers open $\pk = (\vec{\widetilde{a}}, \vec{\widetilde{u}},
\vec{\widetilde{v}}, \vec{\widetilde{w}})$. Each signer keeps $\shr{\sk} =
(\shr{\mat{K}}, \shr{\vec{b}}, \shr{\vec{x}}, \shr{\vec{y}})$.

    \item[$\shr{\state} \sample \TSPSoffsign(\shr{\sk})$ :] The signers sample
$\shr{\rho}, \shr{\sigma} \sample \algo{RandShare}()$, and compute
      %
      \begin{gather*}
      %
	\shr{\rho\sigma} \sample \algo{Mul}(\shr{\rho}, \shr{\sigma}), \\
      %
	\shr{\vec{r}} = \shr{\vec{b}^{\,\rho}}, \ \shr{\vec{s}} =
\shr{\vec{b}^{\,\rho\sigma}}, \ \shr{\vec{x}^{\,\rho}}, \
\shr{\vec{y}^{\,\rho\sigma}} \sample \algo{ExpGrp}_1,
      %
      \end{gather*}
      %
      the second line as eight calls to $\algo{ExpGrp}_1$, one per entry, in one
round.
      Each signer then computes locally $\shr{\widetilde{q}} =
\ggen^{\shr{\sigma}}$ and $\shr{\vec{z}^*} = \shr{\vec{x}^{\,\rho}} \cdot
\shr{\vec{y}^{\,\rho\sigma}}$, and stores $\shr{\state} = (\shr{\widetilde{q}},
\shr{\vec{r}}, \shr{\vec{s}}, \shr{\vec{z}^*})$.

    \item[$\SPSsig_i \leftarrow \TSPSonsign(\sk_i, \state_i, \vec{\msg})$ :]
Each signer computes
      %
      \begin{equation*}
      %
	\shr{\vec{z}} = \shr{\vec{z}^*} \cdot (\gen, \vec{\msg})^{\shr{\mat{K}}},
      %
      \end{equation*}
      %
      and outputs
$(\shr{\widetilde{q}}, \shr{\vec{r}}, \shr{\vec{s}}, \shr{\vec{z}})$.

    \item[$\SPSsig \leftarrow \TSPScombine(\{\SPSsig_i\}_{i \in Q})$ :]
Reconstruct every group element from the shares of the $t+1$ signers in $Q$, and
output $\SPSsig = (\widetilde{q}, \vec{r}, \vec{s}, \vec{z})$.

    \item[$\{0,1\} \leftarrow \TSPSverify(\pk, \vec{\msg}, \SPSsig)$ :] Output
$1$ if and only if
      %
      \begin{equation*}
      %
	\pairing(\vec{z}, \vec{\widetilde{a}}) = \pairing((\gen,
\vec{\msg}), \vec{\widetilde{u}})\, \pairing(\vec{r},\vec{\widetilde{v}})\,
\pairing(\vec{s}, \vec{\widetilde{w}}) \ \wedge \ \pairing(r_1, \widetilde{q})
= \pairing(s_1, \ggen).
      %
      \end{equation*}

  \end{description}

\end{protocolbox}

\begin{theorem}\label{thm:tsps-security} Let the MPC protocol be private and
secure with abort against up to $t$ adaptively corrupted signers, over a linear
$(t+1, n)$ threshold sharing. Then
the threshold $\AGHO$ and $\KPW$ schemes are correct. The threshold $\AGHO$
scheme is adp-TS-UF-1 secure in the GGM, and the threshold $\KPW$ scheme is
adp-TS-UF-1 secure under SXDH, in both
cases with the same loss as the generic construction.  \end{theorem}
